\section{YOLO11-DLA}
\label{sec:method}
\begin{figure}[t]
\centering
\resizebox{\columnwidth}{!}{%
\begin{tikzpicture}[
    node distance=2.6mm and 3.2mm,
    box/.style={draw, rounded corners=2pt, align=center, minimum height=7.5mm,
                inner xsep=3pt, font=\scriptsize},
    kept/.style={box, fill=green!10},
    sub/.style={box, fill=blue!12, very thick},
    io/.style={box, fill=gray!12},
    host/.style={box, fill=orange!16},
    ar/.style={-{Latex[length=1.6mm]}, thick}
]
\node[io] (img) {image\\$1\!\times\!3\!\times\!640^2$};
\node[kept, right=of img] (bb) {preserved YOLO11\\stem, C3k2, SPPF};
\node[sub, right=of bb] (mix)
  {\textbf{S1}: C2PSA\\$\rightarrow$ C2DLA};
\node[kept, below=5.5mm of bb] (neck) {preserved PAN\\neck (P3--P5)};
\node[sub, right=of neck] (head) {\textbf{S2}: Detect\\$\rightarrow$ DetectDLA};
\node[io, right=of head] (maps) {\texttt{pred\_p3/4/5}\\$144\!\times\!\{80,40,20\}^2$};
\node[host, below=5.5mm of head] (post)
  {host task logic: sparse DFL, sigmoid, NMS, box scaling};

\draw[ar] (img) -- (bb);
\draw[ar] (bb) -- (mix);
\draw[ar] (mix.south) -- ++(0,-2.7mm) -| (neck.north);
\draw[ar] (neck) -- (head);
\draw[ar] (head) -- (maps);
\draw[ar] (maps.south) -- ++(0,-2.7mm) -| (post.north);

\node[draw=blue!65, dashed, rounded corners,
      fit=(img)(bb)(mix)(neck)(head)(maps), inner sep=3.2pt] (engine) {};
\node[font=\scriptsize\bfseries, blue!60!black, above=0.4mm of engine.north west,
      anchor=west] {DLA engine: one loadable, all tensors 4-D, static shapes};
\end{tikzpicture}%
}
\caption{Accelerator boundary and the two substitutions. The learned graph inside the
dashed region compiles to one DLA loadable in the evaluated configuration.
Preprocessing, physical-layout conversion, and detection decoding run on the host. S1 removes the
attention matrix multiplications, S2 removes the rank-three decoded output
interface. Green blocks are inherited from the reference model unchanged.}
\label{fig:boundary}
\end{figure}

\subsection{Design Objective and Scope}
Let $G_0$ denote the reference learned graph and $G$ the adapted graph. We
seek a useful detector subject to
\begin{equation}
 \operatorname{Build}_{\mathrm{DLA}}(G;\theta,\mathrm{fallback}=0)
 = E, \qquad N_{\mathrm{GPU}}(E)=0,
 \label{eq:accept}
\end{equation}
where $\theta$ fixes the target, compiler version, precision, shapes, and
I/O formats, $E$ is the serialized engine, and $N_{\mathrm{GPU}}$ counts
GPU-assigned learned layers after optimization. For the evaluated artifact we
also require one DLA loadable. These are deployment criteria, not terms in the
training loss. Accuracy and application performance must be evaluated after
this feasibility gate.

We retain the convolutional stem, C3k2 stages, SPPF module, feature-pyramid
connections, and detection towers of YOLO11n. Two module substitutions define
the adaptation:
\begin{equation}
 \mathrm{C2PSA}\longrightarrow\mathrm{C2DLA},\qquad
 \mathrm{Detect}\longrightarrow\mathrm{DetectDLA}.
 \label{eq:substitutions}
\end{equation}
The first changes a learned computation; the second changes the export
boundary while preserving the ordinary training path. This separation permits
an architecture comparison on the GPU and a placement comparison using the
same adapted graph. It does not make the stock and adapted engine interfaces
identical.

\subsection{C2DLA: Channel Mixing with Local Context}
\label{sec:c2dla}
\begin{figure}[t]
\centering
\resizebox{\columnwidth}{!}{%
\begin{tikzpicture}[
    op/.style={draw, rounded corners=2pt, align=center, minimum height=6.4mm,
               minimum width=8.5mm, inner xsep=2.5pt, font=\scriptsize},
    conv/.style={op, fill=blue!10},
    act/.style={op, fill=violet!12},
    io/.style={op, fill=gray!14, minimum width=6mm},
    sum/.style={draw, circle, inner sep=0.5pt, font=\scriptsize},
    ar/.style={-{Latex[length=1.6mm]}, thick}
]
\node[io]   at (0.25,-0.65) (b)   {$X$};
\node[conv] at (1.45,0)     (pk)  {$\phi_k$};
\node[act]  at (2.75,0)     (sm)  {channel\\softmax};
\node[conv] at (4.05,0)     (pa)  {$\phi_a$};
\node[conv] at (1.45,-1.30) (pv)  {$\phi_v$};
\node[conv] at (2.75,-1.30) (dw)  {DW $3\!\times\!3$};
\node[sum]  at (4.95,-0.65) (s1)  {$+$};
\node[conv] at (5.75,-0.65) (po)  {$\phi_o$};
\node[sum]  at (6.55,-0.65) (r1)  {$+$};
\node[conv] at (7.40,-0.65) (ffn) {FFN};
\node[sum]  at (8.25,-0.65) (r2)  {$+$};
\node[io]   at (9.00,-0.65) (out) {$X''$};

\draw[ar] (b.east) -- ++(0.18,0) |- (pk.west);
\draw[ar] (b.east) -- ++(0.18,0) |- (pv.west);
\draw[ar] (pk) -- (sm);
\draw[ar] (sm) -- (pa);
\draw[ar] (pv) -- (dw);
\draw[ar] (pa.east) -| (s1.north);
\draw[ar] (dw.east) -| (s1.south);
\draw[ar] (s1) -- (po);
\draw[ar] (po) -- (r1);
\draw[ar] (r1) -- (ffn);
\draw[ar] (ffn) -- (r2);
\draw[ar] (r2) -- (out);
% residual around the mixer, Eq. (8)
\draw[ar] (0.25,-0.34) -- (0.25,0.80) -- (6.55,0.80) -- (r1.north);
% residual around the FFN, Eq. (9)
\draw[ar] (6.55,-0.96) -- (6.55,-1.95) -- (8.25,-1.95) -- (r2.south);
\end{tikzpicture}%
}
\caption{The inner C2DLA block. Channel softmax makes the upper branch
input-dependent at each position; the $3\times3$ depthwise branch supplies local
spatial context. Both residuals, the feed-forward block, and the enclosing
split--process--merge structure are retained. The replacement does not compute
global spatial attention.}
\label{fig:c2dla}
\end{figure}

C2DLA keeps C2PSA's split--process--merge structure. An initial $1\times1$
projection produces two channel groups; one bypasses the inner block and the
other is processed before concatenation and a final projection. If $X$ is the
processed group with $C$ channels, define
\begin{align}
 K &= \phi_k(X), \qquad V=\phi_v(X), \label{eq:kv}\\
 A_{bkhw} &= \frac{\exp K_{bkhw}}{\sum_{j=1}^{d}\exp K_{bjhw}},
 \label{eq:chansoftmax}\\
 Z &= \phi_o\bigl(\phi_a(A)+\psi_{3\times3}(V)\bigr),
 \label{eq:mix}\\
 X' &= X+Z, \label{eq:res1}\\
 X'' &= X'+\operatorname{FFN}(X'). \label{eq:res2}
\end{align}
Here $d=\max(\lfloor C/4\rfloor,8)$; $\phi_k:C\to d$,
$\phi_v:C\to C$, $\phi_a:d\to C$, and $\phi_o:C\to C$ are
$1\times1$ projections. The operator $\psi$ is a $3\times3$ depthwise
convolution. The FFN expands channels from $C$ to $2C$, applies SiLU, and
projects back to $C$. Batch normalization accompanies the relevant convolution
modules and is folded into them for inference; $\phi_a$ is a bias-free
convolution without batch normalization.

For the nano configuration, the enclosing block receives 256 channels, the
processed group has $C=128$, the bottleneck has $d=32$, and one inner block
operates at the $20\times20$ scale. All operations preserve the four-dimensional
batch--channel--height--width representation. Unlike the original attention,
there is no spatial affinity matrix, matrix multiplication, head transpose,
or normalization over image positions.

The softmax output remains input-dependent: changing a pixel's feature vector
changes its channel weights. However, this branch independently transforms
each spatial position, and its projection back to $C$ channels does not provide
global information exchange. The depthwise branch supplies local spatial
context. The substitution therefore changes the hypothesis class and requires
empirical accuracy evaluation; it is not an algebraic equivalence.

The key and value projections are separate convolutions. This avoids introducing
a fused-QKV extraction subgraph in the replacement branch. Static channel
splits elsewhere in the retained network remain present and compile successfully
in the measured configuration. Thus, the implementation does not depend on a
claim that every split or SiLU is unsupported on DLA.

\subsection{DetectDLA: A Static Prediction Boundary}
\label{sec:boundary}
For each pyramid level $\ell\in\{3,4,5\}$, the retained regression and
classification towers produce $R_\ell$ and $S_\ell$. DetectDLA exports
\begin{equation}
 P_\ell=\operatorname{Concat}_{C}(R_\ell,S_\ell)
 \in\mathbb{R}^{B\times(4r+n_c)\times H_\ell\times W_\ell},
 \label{eq:packed}
\end{equation}
with $r=16$ regression bins and $n_c=80$ COCO classes. For batch one and
$640\times640$ images, the outputs are $1\times144\times80\times80$,
$1\times144\times40\times40$, and $1\times144\times20\times20$ at strides
8, 16, and 32. They are named \texttt{pred\_p3}, \texttt{pred\_p4}, and
\texttt{pred\_p5}. Classification values are logits; box channels are
unnormalized distributions, not decoded coordinates.

In export mode, the head stops before flattening, cross-scale concatenation,
DFL expectation, grid decoding, and sigmoid. NMS is also outside the engine.
During training and ordinary PyTorch inference, DetectDLA invokes the inherited
Detect path, retaining its losses and assignment behavior. No direct-regression
or one-to-one head is introduced in this study.

The static maps make the execution boundary explicit. Stock YOLO11n's engine
instead emits a decoded $1\times84\times8400$ tensor. At FP16, the adapted
outputs contain 2,419,200 bytes, compared with 1,411,200 bytes for the stock
output, a factor of 1.714. Consequently, a lower raw-engine latency does not
by itself establish a faster complete detector: the adapted pipeline transfers
more data and performs additional host decoding.

\begin{table}[t]
\caption{The evaluated adaptation. Unchanged denotes retained structure, not
identical trained weights.}
\label{tab:design}\centering\footnotesize
\setlength{\tabcolsep}{3pt}
\begin{tabular}{p{.26\columnwidth}p{.31\columnwidth}p{.32\columnwidth}}
\toprule
Component & YOLO11n & \model \\
\midrule
Stem, C3k2, SPPF, neck & Reference structure & Unchanged \\
Terminal mixer & Spatial self-attention & Channel softmax and local DWConv \\
Box/class towers & Distributional boxes; class logits & Unchanged \\
Training head & Original Detect & Inherited Detect path \\
Export output & One decoded rank-3 tensor & Three raw rank-4 tensors \\
Host detection logic & NMS & DFL, sigmoid, NMS \\
\bottomrule
\end{tabular}
\end{table}

\subsection{Training and Initialization}
The compatible structure permits transfer of shape-matched parameters from a
YOLO11 checkpoint before fine-tuning. The accompanying code also contains an
optional block-aware initializer that gathers compatible key/value projection
channels and copies positional and output projections. The new channel
projection has no direct self-attention counterpart. This optional path and a
staged freeze/unfreeze recipe are implementation capabilities; the retained
checkpoint does not establish that either was used for its training run.

The checkpoint actually evaluated here contains a 500-epoch training history,
with batch size 32, image size 640, seed 0, automatic optimizer selection,
pretrained initialization requested, AMP enabled, and no frozen layers in the
stored training arguments. It contains 2,615,696 unfused parameters. The exact
initialization source and realized optimizer settings must be recovered before
attributing accuracy to a particular transfer strategy.
\missing{P2}{archive the original training command, initialization checkpoint,
optimizer state/log, dataset identity, and independently repeated validation}.
